\documentclass{article}
\usepackage[utf8]{inputenc}
\usepackage{pifont}	% Postscript ZapfDingbats
% bibliography
\usepackage[backend=biber, style=alphabetic]{biblatex}
\addbibresource{cybergoth.bib}
\begin{document}

\paragraph{Introduction} 
Hello, this is a short presentation on reading Crash through Plato's The Republic 
and part of a wip, that is exploring the CyberGothic.

\paragraph{The Republic}
Both Crash and The Republic must be read with some caution;
The Republic, written long before the liberalism of the enlightenment, 
appears at the beginning as entirely redundant; 
we are tirelessy lulled into human betterment and
might want to skip the scan of a lenghtly brainwash, 
but do carry on as the morality play swiftly 
takes form of a Greek tragedy and you suddenly wake in the totalitarian polis 
of Catherine and James Ballard. 
Plato is minecrafting the ideal state from the bottom up, 
starting with a small community of artisan farmers;
however soon the exchange of basic needs does not meet the 
high standards of a European democracy anymore 
and we come to see specialization of atomic scale.
At this point the relations to what matters most 
are elongated to such extent 
that the "language of everything" (Artaud) is no longer intelligible.

\paragraph{the milieu of the cave} 
In tradition of Gothic horror and 
in possession of the techn\={e} of tunneling,
J. G. Ballard is producing knowledge out of psuch\={e} (\ding{38}schuke); 
The analogy of the car accident serves to demonstrate 
possible effects of sensorial shutdown and sudden exposure to darkness, 
a hurtling from convertible space 
into the derelict depths of immutable time, 
haunted by a scarred and severly wounded soul.
Vaughan is repelled by its own simulacra, 
which are deformed past all recognition 
and display morbid behaviour,
and yet do not want to refrain from driving.
In consequence "The Good as Ultimate Object of Knowledge" (rep, 2007, 226) 
plots the rational act of murder 
and in time of death the transparency of the Real is splashing
onto the windshield, 
which in-itself is opaque protective matter 
but porous in letting pass any arbitrary concept while its 
form is positioned towards the sun. 
The exploit shatters the defective Good 
into a multiplicity of undefined prisms,
which are keys for possibilities of new perspectives 
and in a cascade of confetti welcome the reborn.

\paragraph{Mathematical reasoning (dianoia) in Ballard}
There is no doubt that Ballard 
utilizes the principle of the Good and strives for perfection, 
but his ideal writing proves all the more astonishing 
as Dianoia is frequently illustrated in his work. 
Many, if not all of his protagonists possess the extraordinary ability to 
detect a pleasant symmetry between the dreamed geometries of 
consciousness (architecture) and the anatomies of deformed bodies, which gives profound reason to accept, 
that both are the eerie structures of subconscious design. 
A truth to which baffled readers may not be privy 
and in consequence dismiss dianoia as weird visions of the eccentric, or even worse, 
equate frequent blurring of boundaries between life and unlife 
with sharp superimpositions of dead figures. 
Such infertile technology of informal machine learning furthers 
chronic progression of visible shadows in infirmary D, 
where C is observed and validated through temporal form and it 
must be held against the sun that diagnosis by medium of contrast is 
deduced from subjacent stages of the visible (hor\={a}ton). 
To gain traction on life of sector C, 
which in fact remains immune to the strict isolation policy of a sclerotic Peras, 
it might be wise to keep the corpse uncovered 
and activate cooling function A (present /future) 
that in combination with Dianoia B 
allows for apyretic moving through forms to forms. 
Mathematical Formalism locates the transforming body 
across a maze of permanent geometries, 
in which in fit of self-actualisation is parading self-inflicted wounds 
for a wilful contamination with time to partake in the festivities 
of the NOW and all its infectious possiblities.
In advanced stages of Dianoia, i.e. 
seizing the world as governing empire of interrelated signs, 
psuch\={e} (\ding{38}schuke) contracts discrete form, a condition which 
Hilbert diagnoses as
"the intuitive a priori mode of thought"(Tasi\'{c} 2001, 68).

\paragraph{Form in CyberGothic} 
J. G. Ballard provokes thought with radical empiricism - 
steering eyes straight towards body fluids, genitalia, 
and senseless violence, 
full in line with the abysall terrors of petroculture, 
partly visible in the occidental mise en sc\`{e}ne 
of car parks, flyovers and underpasses.
If Crash is not acknowledged for exceeding literary excellence 
on account of Ballard's urge to abstraction (Abstraktionsdrang), 
it definitely generates likes for profit 
by satisfying the urge to empathy (Einf\"{u}hlungsdrang).
The two drives produce the respective 
aesthetics of Kunstsch\"{o}n and Natursch\"{o}n; 
Worringer emphasises the importance of 
keeping the two lines apart, which are 
keys to a precise analysis of 
the psychological conditions of a society 
and urges to take a closer look at the forms of cultural products 
instead of being absorbed in the
absurd but lucrative evaluation of pathological consumer behaviour.(wor, 2)
One better takes fright in a scenario 
where Natursch\"{o}n equals Kunstsch\"{o}n 
as Thanatos sits behind the steering wheel 
and encourages fascist simpleton 
to live out its passive sadistic voyeurism and 
find pleasure in 
a mindless bore of nature, or act of empathy (Einf\"{u}hlungsakt). 
The death drive gets even more intense when all effort and life energy
is turned outwards in awe of representation, (wor, 4) 
and indeed linguistic turns of a Neural Network can be quite impressive 
since it claims deep inquiries into the dimension of time 
and literally traces the organic aesthetic 
for smearing the human construct upon the physical world, 
which from the pole of spiritual perspective is both abstraction and shadow. 
Worringer insists on the inner "will to form" 
as precondition of apperceptive activity, 
expressed in materialisation WITHIN the object.(wor, 8)

\nocite{*}	%prints all .bib entries \printbibliography%[title={Bibliography}] 
\end{document}
